% 中文摘要頁
\begin{ZhAbstract}
    \begin{ZhAbstractItems}
        \noindent \text 關鍵詞：胸部 CT、COPD、量化生物標記、肺氣腫、自動化分類

    \end{ZhAbstractItems}

    \begin{ZhAbstractDescription}
        本研究提出一套以胸部 CT 影像為基礎之 COPD 輔助分析流程，從肺部分割、定量特徵擷取到分類模型訓練，建立可解釋之自動化判讀系統。系統以 54 筆臨床 CT 影像為資料來源，擷取 12 項具臨床意義之量化生物標記，包括全肺與分葉肺氣腫比例、血管與氣道相關指標、總肺容積及肺動脈直徑，並以五折交叉驗證訓練全連接神經網路進行 Normal/Abnormal 分類。實驗結果顯示，模型可達 92.59\% 之整體準確率與 0.9784 之 AUC；單一特徵分析亦顯示肺氣腫比例與肺容積對分類最具區辨能力。

        除量化生物標記外，本研究亦進一步將 3D CT 影像直接輸入深度模型，探討 COPD 相關肺功能表徵之預測與分級。實驗以 Mamba 與 Hybrid Mamba-Attention 架構為核心，分別進行肺功能塌陷角度回歸、角度三分類、排除中間灰區之極端二分類，以及 GOLD 嚴重度四分類。結果顯示，Hybrid Mamba-Attention 於角度回歸可達 15.85 度之 MAE 與 0.6699 之 Pearson 相關係數；在角度三分類與 GOLD 四分類中，結合資料增強與平衡取樣後，分別達到 0.5713 與 0.5323 之 Macro-F1，Balanced Accuracy 則為 0.6319 與 0.5222。於排除 132 至 151 度灰區後之極端二分類任務中，模型可達 0.8177 之 Macro-F1 與 0.8289 之 Balanced Accuracy，顯示合理的標籤設計與類別不平衡處理有助於提升小樣本 COPD CT 分類之可靠性。整體而言，胸部 CT 定量生物標記與 3D 深度影像模型可相互補足，作為 COPD 輔助診斷、肺功能表徵預測與疾病嚴重度評估之有效依據。
    \end{ZhAbstractDescription}
    
\end{ZhAbstract}
